\chapter{Introduction}

\section{Background}
SSH is the de-facto standard for remote shell access on Linux. Adversaries
brute-force the password authentication by automating tens of thousands of
``Accepted password`` / ``Failed password`` attempts until a valid
credential is guessed. Distributed brute-force campaigns make the threat
hard to mitigate with naïve static rules.

\section{Problem statement}
Pure static rules (e.g. ``> N failed attempts in T seconds``) produce both
false negatives when attackers throttle their attacks and false positives
when a legitimate user mistypes a password. Adaptive thresholding mitigates
some of these errors but it does not exploit richer contextual signals
(unique usernames, port diversity, time-of-day). This project investigates
whether fusing rules with machine-learning classifiers gives a meaningful
improvement.

\section{Objectives}
\begin{enumerate}
  \item Re-implement the Capstone 1 rule engine (static, adaptive, hybrid rules) from its published abstract.
  \item Engineer per-IP features and train supervised (Random Forest, XGBoost) and unsupervised (Isolation Forest) models.
  \item Build a hybrid decision layer combining rule scores and ML probabilities.
  \item Evaluate accuracy, F1, false-positive rate and latency versus the rule-only baselines.
  \item Package a reproducible pipeline and demo dashboard.
\end{enumerate}

\section{Scope}
The project targets Linux ``auth.log``/``secure`` files. Cross-platform
analysis (Windows Event Log, etc.) is left as future work per the abstract.

\section{Report outline}
Chapter 2 reviews related work. Chapter 3 details the proposed hybrid
architecture. Chapter 4 documents the implementation. Chapter 5 reports
results. Chapter 6 discusses findings and limitations. Chapter 7 concludes.
